% Part: turing-machines
% Chapter: undecidability
% Section: enumerating-tms

\documentclass[../../../include/open-logic-section]{subfiles}

\begin{document}

\olfileid{tur}{und}{enu}
\olsection{ٹیورنگ مشیناں دی گنتی وار فہرست}

\begin{explain}
اسی دکھا سکدے آں کہ ساریاں ٹیورنگ مشیناں دا مجموعہ !!{enumerable}
اے۔ ایہ ایس گل توں نکلدا اے کہ ہر ٹیورنگ مشین دا محدود بیان
دتا جا سکدا اے۔ حالتاں دا مجموعہ تے پٹی دے نشاناں دا مجموعہ
دونوں محدود نیں۔ حالتی تبدیلی دا فنکشن $Q \times \Sigma$
توں $Q \times \Sigma \times \{\TMleft, \TMright, \TMstay\}$
تک اک جزوی فنکشن اے؛ ایس لئی جن محدود دلیلی جوڑیاں اُتے
ایہ تعریف شدہ اے، اوہناں لئی ایس دیاں قدراں فہرست وچ لکھ
کے ایہنوں وی متعین کیتا جا سکدا اے۔

ایہ گل اپنی حد تک ٹھیک اے، پر اتھے اک باریک فرق اے۔ ٹیورنگ
مشین دی تعریف نے ایس گل اُتے کوئی پابندی نہیں لائی کہ حالتاں
دے مجموعے تے پٹی دے حروف دے !!{element}s کیہڑے ہو سکدے نیں۔
مثلاً ہر حقیقی عدد لئی، تکنیکی طور اُتے، اک ایسی ٹیورنگ مشین
اے جیڑی اُس عدد نوں اپنی حالت بنا لَیندی اے۔ پر مشین دا
\emph{عمل} ایس گل توں آزاد اے کہ حالت تے حرف دے طور اُتے کیہڑیاں
اشیا ورتیاں گئیاں نیں۔ \olref{fig:variants} وچ دتیاں دو ٹیورنگ
مشیناں ویکھو۔
\begin{figure}
\begin{center}
\begin{tikzpicture}[->,>=stealth',shorten >=1pt,auto,node distance=2.8cm,
                    semithick]
  \tikzstyle{every state}=[fill=none,draw=black,text=black]

  \node[initial,state]         (A)              {$q_0$};
  \node[state]         (B) [right of=A] {$q_1$};

  \path (A) edge [bend left] node {\TMtrans{\TMstroke}{\TMstroke}{\TMright}} (B)
        (B) edge [loop above] node {\TMtrans{\TMblank}{\TMblank}{\TMright}} (B)
            edge [bend left] node {\TMtrans{\TMstroke}{\TMstroke}{\TMright}} (A);
\end{tikzpicture}\\
\begin{tikzpicture}[->,>=stealth',shorten >=1pt,auto,node distance=2.8cm,
  semithick]
\tikzstyle{every state}=[fill=none,draw=black,text=black]

\node[initial,state]         (A)              {$s$};
\node[state]         (B) [right of=A] {$h$};

\path (A) edge [bend left] node {\TMtrans{A}{A}{\TMright}} (B)
(B) edge [loop above] node {\TMtrans{\TMblank}{\TMblank}{\TMright}} (B)
edge [bend left] node {\TMtrans{A}{A}{\TMright}} (A);
\end{tikzpicture}
\end{center}
\caption{\emph{جُفت} مشین دے وکھرے روپ}
\ollabel{fig:variants}
\end{figure}
ایہ دو خاکے دو مشیناں نوں ظاہر کردے نیں: $M$ دی پٹی دا
حرفی مجموعہ $\Sigma = \{\TMendtape,\TMblank,\TMstroke\}$
تے حالتاں دا مجموعہ $\{q_0,q_1\}$ اے؛ $M'$ دا حرفی
مجموعہ $\Sigma' =
\{\TMendtape,\TMblank,A\}$ تے حالات $\{s,h\}$ نیں۔ پر
باقی ہدایات اکّو جیہیاں نیں: $M$، $n$ عدد $\TMstroke$
نشاناں دی لڑی اُتے اُدوں تے صرف اُدوں رُکدی اے جدوں $n$ جُفت
ہووے؛ تے $M'$، $n$ عدد $A$ نشاناں دی لڑی اُتے، اُدوں تے
صرف اُدوں رُکدی اے جدوں $n$ جُفت ہووے۔ اسی صرف $\TMstroke$
دا ناں بدل کے~$A$، $q_0$ دا بدل کے~$s$، تے $q_1$ دا
بدل کے~$h$ رکھیا اے۔ ایہ مثال عام اصول دِندی اے: جے
اک مشین دوجی توں حالتاں تے نشاناں دے ناں ایس طرح بدلݨ نال
ملدی اے، تاں اسی اوہناں نوں اکّو جیہی سمجھ سکدے آں۔
بلکہ اسی مشین دے نشاناں تے حالتاں نوں بس مثبت صحیح اعداد
ہی سمجھ سکدے آں: $\sigma_0$ دی تھاں~$1$، $\sigma_1$
دی تھاں~$2$، وغیرہ؛ $\TMendtape$ نوں~$1$، $\TMblank$
نوں~$2$، وغیرہ۔ ایس طرح \emph{جُفت} مشین
\olref{fig:standard-even} وچ وکھائی مشین بن جاندی اے۔
\begin{figure}
\[\begin{tikzpicture}[->,>=stealth',shorten >=1pt,auto,node distance=2.8cm,
  semithick]
\tikzstyle{every state}=[fill=none,draw=black,text=black]

\node[initial,state]         (A)              {$1$};
\node[state]         (B) [right of=A] {$2$};

\path (A) edge [bend left] node {\TMtrans{3}{3}{\TMright}} (B)
(B) edge [loop above] node {\TMtrans{2}{2}{\TMright}} (B)
edge [bend left] node {\TMtrans{3}{3}{\TMright}} (A);
\end{tikzpicture}
\]
\caption{اک \emph{معیاری} جُفت مشین}
\ollabel{fig:standard-even}
\end{figure}
جس ٹیورنگ مشین دیاں حالتاں تے نشان مثبت صحیح اعداد ہون،
اسی اوہنوں \emph{معیاری} مشین کہہ سکدے آں تے ہن توں صرف
ایہناں مشیناں اُتے غور کر سکدے آں۔\footnote{“معیاری
مشین” اک رائج معیاری اصطلاح نہیں۔}

ساڈا مقصد ایہ دکھانا سی کہ ٹیورنگ مشیناں دا مجموعہ
!!{enumerable} اے۔ اُپرلی گل دے پیشِ نظر ایس لئی ایہ
دکھانا کافی اے کہ معیاری ٹیورنگ مشیناں دا مجموعہ
!!{enumerable} اے۔ فرض کرو ساڈے کول اک معیاری مشین
$M = \tuple{Q, \Sigma, q_0, \delta}$ اے۔ اسی مثبت صحیح
اعداد دی اک محدود لڑی نال ایہدا بیان کِویں دیواں؟ پہلے
حالتاں دی تعداد، فیر خود حالات، نشاناں دی تعداد، خود نشان
تے ابتدائی حالت فہرست وچ لکھانگے۔ (چیتے رکھو کہ ایہ
سارے مثبت صحیح اعداد نیں، کیوں جے $M$ معیاری مشین اے۔)
ہن~$\delta$ دا کیہ کریے؟ ممکنہ آرگومنٹاں، یعنی
$\tuple{q,\sigma}$ جوڑیاں، دا مجموعہ محدود اے، کیوں جے
$Q$ تے~$\Sigma$ محدود نیں۔ ایس لئی~$\delta$ دی ساری
معلومات بس اوہناں $5$-رُکنی سلسلیاں
$\tuple{q, \sigma, q', \sigma', d}$ دی محدود فہرست اے
جنہاں لئی $\delta(q,\sigma) = \tuple{q', \sigma', D}$
ہووے۔ اتھے $d$ اک عدد اے جیڑا سمت~$D$ نوں کوڈ کردا
اے (مثلاً $1$ نال~$\TMleft$، $2$ نال~$\TMright$، تے
$3$ نال~$\TMstay$)۔

ایس طرح ہر معیاری ٹیورنگ مشین دا بیان مثبت صحیح اعداد دی
اک محدود فہرست، یعنی اک سلسلے $s_M \in
(\PosInt)^*$، نال دتا جا سکدا اے۔ مثلاً معیاری
\emph{جُفت} مشین نوں ایہ سلسلہ کوڈ کردا اے:
\[
2, \underbrace{1, 2}_Q, 3, \overbrace{1, 2, 3}^\Sigma, 1, \underbrace{1, 3, 2, 3, 2}_{\delta(1,3) = \tuple{2,3,R}}, 
\overbrace{2, 2, 2, 2, 2}^{\delta(2,2) = \tuple{2,2,R}},
\underbrace{2, 3, 1, 3, 2}_{\delta(2,3) = \tuple{1,3,R}}.
\]
\end{explain}

\begin{thm}
$\Nat$ توں~$\Nat$ تک ایسے فنکشن موجود نیں جیڑے ٹیورنگ
مشین نال حساب نہیں ہو سکدے۔
\end{thm}

\begin{proof}
سانوں پتا اے کہ مثبت صحیح اعداد دے محدود سلسلیاں
~$(\PosInt)^*$ دا مجموعہ !!{enumerable} اے
(\cref{sfr:siz:zigzag:prob:posint-star})۔ ایس توں نکلدا
اے کہ معیاری ٹیورنگ مشیناں دے بیان، جو~$(\PosInt)^*$
دا ذیلی مجموعہ نیں، خود وی گنتی وار فہرست بنن جوگ نیں۔
ہر ٹیورنگ-حسابی فنکشن $\Nat$ توں~$\Nat$ تک کسے نہ
کسے (اصل وچ کئی) ٹیورنگ مشین نال حساب ہوندا اے۔ اوس
مشین دیاں حالتاں تے نشانیاں دے ناں مثبت صحیح اعداد
رکھ کے، خاص طور اُتے $\TMendtape$ نوں~$1$، $\TMblank$
نوں~$2$، تے $\TMstroke$ نوں~$3$ رکھ کے، اسی ویکھدے آں
کہ ہر ٹیورنگ-حسابی فنکشن دا حساب اک معیاری ٹیورنگ مشین
وی کردِی اے۔ ایس دا مطلب اے کہ $\Nat$ توں~$\Nat$ تک
سارے ٹیورنگ-حسابی فنکشناں دا مجموعہ وی گنتی وار فہرست
بنن جوگ اے۔

دوجے پاسے، $\Nat$ توں~$\Nat$ تک سارے فنکشناں دا
مجموعہ !!{enumerable} نہیں اے
(\cref{sfr:siz:red:prob:nat-nat})۔ جے سارے فنکشن
کسی نہ کسی ٹیورنگ مشین نال حساب ہو سکدے ہوندے، تاں
اوہناں نوں حساب کرن والیاں مشیناں دے بیان فہرست وچ
لکھ کے اسی سارے فنکشناں دی فہرست بنا سکدے ہوندے۔
ایس لئی کجھ فنکشن ٹیورنگ مشین نال حساب نہیں ہو سکدے۔
\end{proof}

\begin{prob}
  کی تسی ٹیورنگ مشیناں دا ایسا بیان سوچ سکدے او جس وچ
  حالات تے حرفی نشاناں نوں صاف صاف فہرست وچ لکھن
  دی لوڑ نہ ہووے؟ تسی “معیاری” مشین دا اپنا تصور
  تعریف کر سکدے او، پر ایہ وی دسو کہ تہاڈے نویں معنیٰ
  وچ ہر ٹیورنگ مشین دی حسابی کارروائی اک “معیاری”
  مشین کیوں کر سکدی اے۔
\end{prob}

\end{document}
